\section{A determinantal model for every contact order}
\label{sec:determinantal-v167}
The distinction between a curve in a Hilbert scheme and a fibre of
its parameter space is essential. This section constructs the former
Hilbert graph over the whole polynomial coefficient space. Later we
return to flat limits of the fibres of that graph. The two operations
have different Hilbert polynomials and different universal properties.

Fix an integer $a\geq1$ and set
\[
 B_a=\{(f,g,r):f\text{ monic of degree }a,\quad
                         \deg g,\deg r<a\}.
\]
Write $[s:t]$ for the source coordinates and $x=t/s$. Homogenization
to degree $a$ gives forms $F_0,F_1,F_2$. Let $U_a\subset B_a$ be
the open locus where these forms have no common zero. The graph
$C_b\subset X=\Pj^1\times\Pj^2$ of $[F_0:F_1:F_2]$ has, for the
Segre polarization, polynomial
\[
 P_a(l)=(a+1)l+1.
\]
The reduced closure of $(b,[C_b])$, $b\in U_a$, is denoted by
$\Gamma_a\subset B_a\times\Hilb^{P_a}(X)$, as in
Section~\ref{sec:Euclidean-boundary-v162}. It is an integral graph,
not a quotient by projective changes of the target.

Put
\begin{equation}\label{eq:hilbert-degree-v167}
 d=a+1,\qquad m=\frac{d(d-1)}2+1,\qquad
 q=dm+1,\qquad N=(m+1)\binom{m+2}{2}.
\end{equation}
For $0\leq i\leq m$ and $j_0+j_1+j_2=m$, let the corresponding
column of the $q\times N$ matrix $M_a$ be the coefficient vector of
\begin{equation}\label{eq:state-matrix-v167}
 s^{m-i}t^i F_0^{j_0}F_1^{j_1}F_2^{j_2}
                   \quad\text{in }\C[s,t]_{dm}.
\end{equation}
Thus its entries are explicit polynomials in the $3a$ coefficients.
Order rows by the exponent of $t$, and order columns
lexicographically in $(i,j_0,j_1,j_2)$. For each $q$-element column
set $J$, write $\Delta_J=\det(M_a)_J$; identically zero determinants
may be omitted. Let
\begin{equation}\label{eq:hilbert-ideal-v167}
                         \mathfrak a_a=(\Delta_J)_J\subset\OO(B_a).
\end{equation}
The size of this presentation is not a complexity bound. In
particular, enumerating every maximal minor is not the recommended
way to specialize one matrix.

\begin{theorem}[Determinantal Hilbert graph]\label{thm:determinantal-v167}
For every $a\geq1$ there is a canonical isomorphism over $B_a$
\begin{equation}\label{eq:rees-hilbert-v167}
 \Gamma_a\simeq\operatorname{Bl}_{\mathfrak a_a}(B_a)
       =\Proj_{B_a}\bigoplus_{n\geq0}\mathfrak a_a^n.
\end{equation}
Under this isomorphism the rational Pl\"ucker coordinates of the
Hilbert map are $[\Delta_J]_J$. The universal curve on this blow-up
is the flat embedded extension of the graph on $U_a$.
The normalization of $\Gamma_a$ is the normalization of this Rees
model. This assertion does not normalize any scheme fibre.

Let $R$ be a discrete valuation ring containing $\C$, with
uniformizer $\tau$, fraction field $K$, and residue field $k$.
For a coefficient arc $\gamma:\Spec R\to B_a$ with
$\gamma(\Spec K)\subset U_a$, set
\begin{equation}\label{eq:determinantal-order-v167}
 \mu(\gamma)=\min_J\operatorname{ord}_{\tau}\Delta_J(\gamma),
 \qquad
 \ell_J(\gamma)=\left(\tau^{-\mu(\gamma)}
                                \Delta_J(\gamma)\right)\bmod\tau.
\end{equation}
The special embedded curve is the unique Hilbert point whose
Pl\"ucker vector is $[\ell_J(\gamma)]_J$. In particular this vector
specifies its scheme structure, not merely its support or cycle.
Two such arcs centred at the same coefficient point have the same
point of $\Gamma_a$ precisely when these nonzero vectors are
proportional over $k$.
\end{theorem}

\begin{proof}
We first check the degree at which the Hilbert point is taken.
The Gotzmann representation of $dl+1$ has $d$ terms of degree one
and $(d-1)(d-2)/2$ terms of degree zero, since
\[
 \sum_{i=0}^{d-1}\binom{l+1-i}{1}
                       =dl+\frac{d(3-d)}2.
\]
Its total number of terms is $m$ in
\eqref{eq:hilbert-degree-v167}. Gotzmann regularity and the
Grassmannian construction of the Hilbert scheme therefore give a
closed immersion into the Grassmannian of rank-$q$ quotients in
degree $m$; see \cite{Gotzmann167,Grothendieck166}. We use the
ordinary projective Hilbert scheme after the Segre embedding
$X\hookrightarrow\Pj^5$, not an unproved multigraded regularity
bound. The restriction
\[
 H^0(\Pj^5,\OO(m))\longrightarrow H^0(X,\OO(m,m))
\]
is onto. All kernels defining subschemes of $X$ contain its fixed
kernel. Hence this Hilbert immersion factors as a closed immersion
into the Grassmannian of quotients of $H^0(X,\OO(m,m))$, of
dimension $N$.

For $b\in U_a$, evaluation on the graph is exactly
\[
 H^0(X,\OO(m,m))\longrightarrow
 H^0(C_b,\OO(m,m))=H^0(\Pj^1,\OO(dm)).
\]
This map is onto by the same regularity theorem. Its matrix is
$M_a(b)$ and its Pl\"ucker coordinates are its $q$-minors. This
proves that $\mathfrak a_a$ is nonzero and has no base point on
$U_a$.

For an integral affine scheme $\Spec A$ and a nonzero ideal
$(h_0,\ldots,h_n)$, the graph closure of $[h_0:\cdots:h_n]$ is
$\Proj A[(h_0,\ldots,h_n)T]$. Indeed, on the chart indexed by a
nonzero $h_i$, both schemes have ring
$A[h_0/h_i,\ldots,h_n/h_i]\subset\operatorname{Frac}(A)$.
These charts prove the assertion scheme-theoretically, with the
Rees algebra rather than the symmetric algebra. Apply this to
$\mathfrak a_a$. The Hilbert scheme is closed in the Pl\"ucker
space, so the closure of its dense graph is this same blow-up.
Pulling back the universal Hilbert family proves the flat-family
assertion. Normalizing the identical integral schemes proves the
normalization assertion.

Over $R$, at least one $\Delta_J(\gamma)$ is nonzero. Dividing
all coordinates by their common valuation gives a primitive
vector in a finite free $R$-module. Its special projective point
is $[\ell_J]$. The Grassmannian and Hilbert equations hold over
$K$ and hence over $R$, since $R$ has no $\tau$-torsion. The
resulting $R$-point is the unique extension by separatedness and
properness. The Hilbert immersion makes equality of its special
projective coordinates equivalent to equality of the embedded
special curves. Keeping the coefficient point supplies the last
assertion for the retained-coefficient graph.
\end{proof}

\begin{remark}[What the determinant construction supplies]
\label{rem:scope-states-v167}
The existence of a graph blow-up and the Hilbert immersion are
classical. The explicit ideal \eqref{eq:hilbert-ideal-v167} identifies
their input for every $B_a$ and makes the following valuation
statements uniform in $a$. It is not a claim that arbitrary
higher-contact fibres have only finitely many isomorphism classes.
It does not identify points on the normalization solely from their
Hilbert coordinates: distinct normalization points can lie over
one graph point. A generically nonincident DVR arc has a unique
lift to that normalization by properness and the generic
isomorphism. Equality of closed Hilbert points is a statement on
$\Gamma_a$, not an assertion that the lifts coincide.
\end{remark}

\begin{proposition}[A finite determinacy bound for each arc]
\label{prop:jet-v167}
Keep the matrix and coordinates above fixed. If two $R$-valued
coefficient arcs agree modulo $\tau^{\mu(\gamma)+1}$ and the first
is generically in $U_a$, the second is also generically in $U_a$
and their embedded special curves agree. Thus a jet of length
$\mu(\gamma)+1$ suffices for this arc. The assertion concerns a
Hilbert point, not the whole formal family.
\end{proposition}
\begin{proof}
Polynomial evaluation preserves congruences. All maximal minors
are therefore congruent modulo $\tau^{\mu+1}$, and at least one
has the same nonzero coefficient in degree $\mu$. The primitive
vectors in \eqref{eq:determinantal-order-v167} agree. To justify
the generic-locus assertion, a triple with common homogeneous
factor of positive degree has all images in
\eqref{eq:state-matrix-v167} divisible by its $m$th power and
cannot span $\C[s,t]_{dm}$. Thus a nonzero maximal minor forces
absence of a common zero over an algebraic closure of $K$.
\end{proof}

\begin{proposition}[Smith data and specialization without enumerating minors]
\label{prop:smith-v167}
Let the nonzero invariant factors of $M_a(\gamma)$ over $R$ have
valuations $e_1\leq\cdots\leq e_q$. Then
\[
 \mu(\gamma)=\sum_{i=1}^q e_i.
\]
The cokernel of this degree-$m$ evaluation map is
$\bigoplus_i R/(\tau^{e_i})$. Its least annihilating exponent is
$e_q$, and its length is $\mu$. The saturated row lattice of
$M_a(\gamma)$ in $R^N$ has special exterior line $[\ell_J]$.
The kernel of the corresponding rank-$q$ quotient over $k$,
viewed in the Segre coordinate ring, is the degree-$m$ part of the
special ideal. It generates its ideal sheaf.
\end{proposition}
\begin{proof}
Smith normal form follows by choosing an entry of least valuation,
using it to clear its row and column, and repeating. Elementary
invertible operations preserve the ideal of maximal minors. In
normal form that ideal is generated by
$\tau^{e_1+\cdots+e_q}$, which proves the formula and the cokernel
statements. Over a DVR a saturated submodule of a finite free
module is a direct summand. Saturating the row lattice divides
its exterior line by exactly this common factor. Reduction gives
the stated quotient. Gotzmann regularity, applied to the Hilbert
point already constructed, says that its saturated homogeneous
ideal is generated in degrees at most $m$, and its sheaf is
generated by the degree-$m$ part.
\end{proof}
The exponent $e_q$ here is an exponent for a finite module in one
fixed projective degree. It is not a bound for saturation of every
affine presentation of a pulled-back parameter graph.

\section{Finite coefficient-state chambers}
\label{sec:state-fan-v167}
A valuation without its leading coefficient loses information.
The following finite statement keeps that information, including
cancellation. Its use of coefficient weights is different from a
one-parameter torus acting on the fixed target. Classical state
polytopes organize initial ideals under the latter action
\cite{BayerMorrison167}; here the weighted variables are the input
coefficients of \eqref{eq:state-matrix-v167}.

Write these coefficients as $z_1,\ldots,z_{3a}$ and expand the
finite set of nonzero minors as
\[
 \Delta_J(z)=\sum_{\alpha\in E}c_{J\alpha}z^\alpha,
\]
where $E$ is the union of their finite monomial supports. Allow a
specified subset of the $z_i$ to be identically zero, and omit the
corresponding monomials first. For this theorem take $R=k[[\tau]]$ (or $k[\tau]_{(\tau)}$),
where $k/\C$ is any field extension, so the residue coefficients
are actual constants. No coefficient-field choice was required
for the preceding general DVR statements. On the remaining
coordinates consider exact monomial arcs
\begin{equation}\label{eq:monomial-arcs-v167}
 z_i=\xi_i\tau^{w_i},\qquad \xi_i\in k^*,\quad
                              w_i\in\mathbb Z_{>0}.
\end{equation}
The constant monic coefficient of $f$ is not among these variables.

\begin{theorem}[Finite decorated chamber classification]
\label{thm:finite-fan-v167}
For each fixed $a$ and each choice of zero coordinates there is an
explicit finite rational polyhedral fan in the nonnegative weight
orthant, and for each relatively open cone a finite locally closed
partition of the residue torus, with the following properties.
The pairs for which \eqref{eq:monomial-arcs-v167} is generically
in $U_a$ are specified by these partitions. On each such pair of
a cone and a residue stratum, the embedded Hilbert limit is given
by a single displayed vector of polynomial functions of $\xi$,
independent of $w$ in that relatively open cone. These vectors
classify the special embedded curves by projective equality.
They include nonreduced curves and embedded scheme structure.

One such fan is the hyperplane arrangement
\begin{equation}\label{eq:fan-v167}
              (\alpha-\beta)\cdot w=0\qquad(\alpha,\beta\in E)
\end{equation}
intersected with the coordinate orthant and all its faces.
The residue partition and the vectors require only finite
polynomial operations on the coefficients $c_{J\alpha}$.
\end{theorem}
\begin{proof}
On the relative interior of a cone, the weak order of all numbers
$\alpha\cdot w$ is fixed. Let $E_1,\ldots,E_t$ be its ordered
equality classes. For each $J$ and $h$ form
\[
 b_{Jh}(\xi)=\sum_{\alpha\in E_h}c_{J\alpha}\xi^\alpha.
\]
Partition the residue torus by the first index $h(J)$ at which
$b_{Jh}$ is nonzero, allowing $h(J)=\infty$ when all are zero.
Each condition is a finite list of polynomial equalities and one
inequality, so there are finitely many locally closed pieces.
Discard the piece on which all $h(J)$ are infinite. On the others
let $h_* =\min_J h(J)$. The answer is the projective vector
\begin{equation}\label{eq:decorated-vector-v167}
 L_J(\xi)=
 \begin{cases} b_{J,h_*}(\xi),&h(J)=h_*,\\0,&h(J)>h_*.
 \end{cases}
\end{equation}
Indeed substitution into every minor is a finite sum of powers of
$\tau$, with its terms grouped exactly as above. Equalities in a
residue stratum account for every possible cancellation, not just
a generic noncancellation assumption. The least surviving exponent
is the common valuation in \eqref{eq:determinantal-order-v167}.
Theorem~\ref{thm:determinantal-v167} proves both classification and
retention of scheme structure. The same rank argument as in
Proposition~\ref{prop:jet-v167} identifies the discarded set with
the triples having a generic common zero.
\end{proof}

\begin{remark}[Finite classification and its scope]
\label{rem:finite-classification-v167}
For fixed residue coefficients, the theorem gives finitely many
embedded limits as the positive integral weights vary. As the
residue coefficients vary, it gives finitely many algebraic
families, not finitely many curves. The arrangement need not be
minimal and does not assert that every one of its walls changes
the curve. The degree bound is uniform for each $a$, but no
claim of efficient enumeration of its potentially enormous minors
is made. Arbitrary formal arcs are treated by the pointwise jet
bound, not by declaring their leading coefficient monomials to be
exact arcs. In particular, repeated cancellation in a formal arc
cannot be discarded by this theorem.
\end{remark}

The construction gives a practical distinction between two sorts
of audit. An individual polynomial arc can be evaluated and its
row lattice saturated, without computing the whole arrangement.
The finite arrangement is a proof of finite classification of
exact monomial profiles. The next section computes the entire
answer on a nontrivial two-dimensional slice, including its
minimal walls and the geometry on those walls.
